\documentclass[a4paper,oneside,11pt]{book}
\usepackage{graphicx}
\usepackage{titling}
\usepackage{geometry}
\usepackage{listings}
\usepackage{xcolor}
\usepackage{float}
\usepackage{indentfirst}
\usepackage{longtable}
\usepackage[T1]{fontenc}
\usepackage{lmodern}
\usepackage[hyphens]{url}
\usepackage{hyperref}
\def\UrlBreaks{\do\/\do-}

\geometry{left=3cm, right=2.5cm, top=3cm, bottom=3cm}
\emergencystretch=5em

\definecolor{codegreen}{rgb}{0,0.6,0}
\definecolor{codegray}{rgb}{0.5,0.5,0.5}
\definecolor{codepurple}{rgb}{0.58,0,0.82}
\definecolor{backcolour}{rgb}{0.97,0.97,0.95}
\definecolor{codeblue}{rgb}{0.16,0.32,0.75}

\lstdefinelanguage{dockerfile}{
    keywords={FROM,RUN,COPY,WORKDIR,EXPOSE,CMD,ARG,ENV,ADD,ENTRYPOINT,VOLUME,USER,LABEL},
    keywordstyle=\color{codeblue}\bfseries,
    basicstyle=\ttfamily\footnotesize,
    sensitive=true,
    comment=[l]{\#},
    commentstyle=\color{codegray}\itshape,
    stringstyle=\color{codepurple},
}

\lstdefinestyle{codestyle}{
    backgroundcolor=\color{backcolour},
    basicstyle=\ttfamily\footnotesize,
    commentstyle=\color{codegreen},
    keywordstyle=\color{codeblue}\bfseries,
    stringstyle=\color{codepurple},
    numberstyle=\tiny\color{codegray},
    numbers=left,
    numbersep=6pt,
    showspaces=false,
    showstringspaces=false,
    breaklines=true,
    tabsize=2,
    frame=single,
    rulecolor=\color{gray},
    captionpos=b,
}
\lstset{style=codestyle}

\title{Laporan Praktikum \\ Konstruksi dan Evolusi Perangkat Lunak\\
Pertemuan 5 \\ Containerisasi Aplikasi Laravel dengan Docker}
\author{Hisyam Al Rayyan / 24/541582/SV/24926}
\date{\today}

\begin{document}

\begin{titlingpage}
\begin{center}
\vspace{4cm}
\begin{huge}
\textbf{\thetitle} \\
\end{huge}
\vspace{2cm}
\includegraphics[height=8cm]{images/logo_ugm.png}\\
\vspace{4cm}
{\Large
\begin{tabular}{rl}
Nama            & : Hisyam Al Rayyan \\
NIM             & : 24/541582/SV/24926 \\
Kelas           & : A2 \\
Dosen Pengampu  & : Galih Malela Damaraji \\
\end{tabular}
}
\vspace{2cm}
\end{center}
\end{titlingpage}

\tableofcontents

%=== BAB 1 ===
\chapter{Pendahuluan}

\section{Tujuan Praktikum}
\begin{enumerate}
    \item Menulis Dockerfile untuk aplikasi Laravel dengan memperhatikan urutan instruksi agar cache Docker bekerja optimal.
    \item Membuat \texttt{.dockerignore} untuk mengecualikan file yang tidak diperlukan dalam image.
    \item Membuktikan efektivitas Docker layer cache melalui tiga kali build dengan kondisi berbeda.
    \item Menjalankan container dan memverifikasi aplikasi dapat diakses dari browser.
    \item Memastikan endpoint API dapat dipanggil dari luar container menggunakan Postman.
\end{enumerate}

\section{Alat dan Bahan}
\begin{enumerate}
    \item Laptop dengan sistem operasi Windows.
    \item Docker Desktop for Windows.
    \item PHP 8.4 dan Composer.
    \item Framework Laravel 13.
    \item Postman untuk pengujian API.
    \item Repository GitHub \texttt{evolusi-pl-24-541582-SV-24926}.
\end{enumerate}

%=== BAB 2 ===
\chapter{Hasil dan Pembahasan}

\section{Dockerfile Aplikasi Laravel}

Dockerfile ditulis dengan memperhatikan urutan instruksi agar Docker layer cache dapat bekerja secara efektif. \texttt{composer.json} dan \texttt{composer.lock} disalin dan dependensi dipasang \textbf{sebelum} kode aplikasi disalin. Dengan cara ini, layer instalasi dependensi hanya dijalankan ulang ketika file \texttt{composer.lock} berubah.

\begin{lstlisting}[language=dockerfile, caption={Dockerfile aplikasi Laravel}]
FROM php:8.4-cli

# Install system dependencies
RUN apt-get update && apt-get install -y \
    git \
    curl \
    libpng-dev \
    libonig-dev \
    libxml2-dev \
    zip \
    unzip \
    sqlite3 \
    libsqlite3-dev \
    && docker-php-ext-install pdo pdo_sqlite mbstring exif pcntl bcmath \
    && apt-get clean && rm -rf /var/lib/apt/lists/*

# Install Composer
COPY --from=composer:latest /usr/bin/composer /usr/bin/composer

WORKDIR /var/www

# === STEP PENTING UNTUK CACHE ===
# Salin composer files DAHULU sebelum kode aplikasi
COPY composer.json composer.lock ./

# Install dependencies (layer ini hanya rebuild jika composer.lock berubah)
RUN composer install --no-dev --optimize-autoloader --no-scripts --no-interaction

# Baru salin seluruh kode aplikasi
COPY . .

# Setup environment
RUN cp .env.example .env \
    && php artisan key:generate \
    && touch database/database.sqlite \
    && php artisan migrate --force \
    && chmod -R 777 storage bootstrap/cache

EXPOSE 8000

CMD ["php", "artisan", "serve", "--host=0.0.0.0", "--port=8000"]
\end{lstlisting}

% Screenshot 1: Dockerfile
%\begin{figure}[H]
%    \centering
%    \includegraphics[width=0.8\textwidth]{5-1.png}
%    \caption{Isi Dockerfile aplikasi Laravel}
%    \label{fig:5-1}
%\end{figure}

\section{File .dockerignore}

File \texttt{.dockerignore} dibuat untuk mengecualikan file dan direktori yang tidak diperlukan dalam image Docker. Hal ini mengurangi ukuran image dan mempercepat proses build.

\begin{lstlisting}[caption={Isi file .dockerignore}]
node_modules
vendor
.git
.env
.env.*
!.env.example
*.log
storage/logs/*
storage/framework/cache/*
storage/framework/sessions/*
storage/framework/views/*
docker-compose.yml
Dockerfile
.dockerignore
tests/
*.md
\end{lstlisting}

% Screenshot 2: .dockerignore
%\begin{figure}[H]
%    \centering
%    \includegraphics[width=0.8\textwidth]{5-2.png}
%    \caption{Isi file .dockerignore}
%    \label{fig:5-2}
%\end{figure}

\section{Pembuktian Docker Layer Cache}

Docker layer cache dibuktikan melalui tiga kali build:

\begin{enumerate}
    \item \textbf{Build pertama}: Semua layer dibangun dari awal. Durasi paling lama karena tidak ada cache yang tersedia.
    \item \textbf{Build kedua}: Tidak ada perubahan sama sekali. Seluruh layer menggunakan cache sehingga build selesai dalam waktu sangat singkat.
    \item \textbf{Build ketiga}: Satu baris kode aplikasi diubah (tanpa mengubah \texttt{composer.lock}). Layer instalasi dependensi menggunakan cache, hanya layer \texttt{COPY . .} dan seterusnya yang dibangun ulang. Jauh lebih cepat dari build pertama.
\end{enumerate}

Perintah yang dijalankan:

\begin{lstlisting}[language=bash, caption={Perintah tiga kali build Docker}]
# Build pertama (semua dari awal)
docker build -t evolusi-app .

# Build kedua (tanpa perubahan - semua cached)
docker build -t evolusi-app .

# Ubah satu baris kode, lalu build ketiga
docker build -t evolusi-app .
\end{lstlisting}

% Screenshot 3: Build pertama (lama)
%\begin{figure}[H]
%    \centering
%    \includegraphics[width=0.8\textwidth]{5-3.png}
%    \caption{Build pertama: semua layer dibangun dari awal}
%    \label{fig:5-3}
%\end{figure}

% Screenshot 4: Build kedua (instant cached)
%\begin{figure}[H]
%    \centering
%    \includegraphics[width=0.8\textwidth]{5-4.png}
%    \caption{Build kedua: seluruh layer menggunakan cache}
%    \label{fig:5-4}
%\end{figure}

% Screenshot 5: Build ketiga (sebagian cached)
%\begin{figure}[H]
%    \centering
%    \includegraphics[width=0.8\textwidth]{5-5.png}
%    \caption{Build ketiga: layer dependensi cached, hanya kode yang dibangun ulang}
%    \label{fig:5-5}
%\end{figure}

\section{Menjalankan Container}

Container dijalankan dengan perintah berikut:

\begin{lstlisting}[language=bash, caption={Menjalankan container Docker}]
docker run -d -p 8000:8000 --name evolusi-container evolusi-app
docker ps
\end{lstlisting}

Aplikasi kemudian dapat diakses melalui browser pada alamat \texttt{http://localhost:8000}.

% Screenshot 6: Browser + docker ps
%\begin{figure}[H]
%    \centering
%    \includegraphics[width=0.8\textwidth]{5-6.png}
%    \caption{Aplikasi berjalan di browser dan output docker ps}
%    \label{fig:5-6}
%\end{figure}

\section{Pengujian API dengan Postman}

Endpoint \texttt{GET /api/tugas} diuji menggunakan Postman untuk memastikan API dapat dipanggil dari luar container. Pengujian memperlihatkan URL yang dipanggil, status HTTP 200, dan isi respons JSON.

\begin{lstlisting}[caption={Endpoint yang diuji dengan Postman}]
GET http://localhost:8000/api/tugas
\end{lstlisting}

% Screenshot 7: Postman request + response
%\begin{figure}[H]
%    \centering
%    \includegraphics[width=0.8\textwidth]{5-7.png}
%    \caption{Postman: request GET /api/tugas dengan status 200 dan respons JSON}
%    \label{fig:5-7}
%\end{figure}

%=== BAB 3 ===
\chapter{Kesimpulan}

Berdasarkan praktikum yang telah dilakukan, dapat disimpulkan beberapa hal berikut:
\begin{enumerate}
    \item Dockerfile berhasil ditulis dengan urutan instruksi yang optimal: \texttt{composer.json} dan \texttt{composer.lock} disalin dan dependensi dipasang sebelum kode aplikasi, sehingga Docker layer cache dapat bekerja secara efektif.
    \item File \texttt{.dockerignore} berhasil dibuat untuk mengecualikan direktori \texttt{node\_modules}, \texttt{vendor}, \texttt{.git}, dan \texttt{.env}, mengurangi ukuran build context.
    \item Efektivitas Docker layer cache terbukti melalui tiga kali build: build ketiga jauh lebih cepat dari build pertama karena layer instalasi dependensi menggunakan cache yang tersimpan.
    \item Container berhasil dijalankan dan aplikasi Laravel dapat diakses dari browser melalui \texttt{http://localhost:8000}.
    \item Endpoint API \texttt{GET /api/tugas} berhasil diuji dengan Postman dan mengembalikan status HTTP 200 beserta data JSON yang benar.
\end{enumerate}

\clearpage

\begin{thebibliography}{9}
\bibitem{dockerDoc}
Docker. (2026). \textit{Dockerfile Reference}. Diakses dari \url{https://docs.docker.com/reference/dockerfile/}.
\bibitem{dockerBestPractices}
Docker. (2026). \textit{Best practices for writing Dockerfiles}. Diakses dari \url{https://docs.docker.com/develop/develop-images/dockerfile_best-practices/}.
\bibitem{laravelDoc}
Laravel. (2026). \textit{Laravel Documentation}. Diakses dari \url{https://laravel.com/docs}.
\end{thebibliography}

\end{document}
